\small
\begin{longtable}{p{2.7cm}p{3.4cm}c p{6.2cm}}
\toprule
Unidad & Subtema & Estado & Justificación técnica y evidencia \\
\midrule\endfirsthead
\toprule Unidad & Subtema & Estado & Justificación técnica y evidencia \\\midrule\endhead
Análisis amortizado & Secuencias de operaciones & A & El costo se evalúa sobre secuencias completas de inserciones; Sec. 2.2 y CSV.\\
& Método agregado & A & Se suman migraciones de todos los resizes y se normalizan por $N$; Sec. 2.2.\\
& Método contable & P & Los créditos por inserción ofrecen una interpretación de cómo financiar migraciones; no es la prueba principal.\\
& Método potencial & P & Se discute como alternativa; ocupación sola no captura redondeo primo e interacción $(\gamma,\tau)$.\\
& Arreglos dinámicos & A & La reserva geométrica y frontera tiempo--espacio fundamentan el redimensionamiento; Sec. 2.3.\\
& Contadores y operaciones & A & Resizes ocasionales, elementos movidos y costo por operación están instrumentados en \texttt{src/}.\\
\midrule
Probabilidad y esperanza & Probabilidad básica & A & La elección de $(a,b)$ define el espacio aleatorio y los eventos de colisión.\\
& Esperanza matemática & A & Se contrasta $\mathbb{E}[C_{pares}]\le n(n-1)/(2m)$; Sec. 2.1 y Fig. 5.\\
& Variables indicadoras & A & Cada par aporta una indicadora; su suma produce $C_{pares}$ y permite linealidad de esperanza.\\
& Amortizado vs. esperado & A & Se distinguen explícitamente garantía de secuencia, esperanza sobre $h$ y media empírica.\\
\midrule
Algoritmos aleatorizados & Las Vegas & P & La tabla siempre conserva corrección; la función aleatoria afecta costo, no respuesta.\\
& Monte Carlo & N & No se acepta probabilidad de respuesta incorrecta: todas las inserciones se verifican.\\
& Probabilidad de error & N & No existe salida aproximada o errónea que deba amplificarse mediante repetición.\\
& Garantías de ejecución & P & Se comparan corrección determinista, costo amortizado y colisiones esperadas.\\
& Selección aleatoria & P & Se eligen aleatoriamente $(a,b)$ y el orden de tratamientos, no pivotes algorítmicos.\\
& Hashing universal & A & Es el mecanismo aleatorizado central; \texttt{src/universal\_hash.py}.\\
\midrule
P, NP y verificadores & Problemas de decisión & N & La pregunta estima rendimiento multiobjetivo, no decide un lenguaje sí/no.\\
& Certificados & N & Las observaciones no requieren testigo para una instancia ``sí''.\\
& Verificadores & N & No se formula certificado; la corrección se comprueba mediante invariantes polinomiales.\\
& Clase P & P & Inserción y validaciones tienen algoritmos polinomiales; no se estudia completitud.\\
& Clase NP & P & Se explica que pertenencia exige problema de decisión y verificador; esa forma no es central aquí.\\
& Decisión vs. optimización & P & Elegir una política Pareto es optimización multiobjetivo, distinta de una decisión formal.\\
\midrule
Reducciones y NP-completitud & Reducciones polinomiales & N & No se compara dificultad de lenguajes mediante transformaciones de instancias.\\
& Dirección de una reducción & N & No se afirma dureza; invertir una reducción no demostraría la implicación buscada.\\
& Prueba de NP-completitud & N & Faltan deliberadamente lenguaje, pertenencia a NP y reducción desde un problema base.\\
& SAT / 3-SAT & N & No modela la selección de parámetros ni la operación de la tabla.\\
& Clique & N & Las colisiones forman grupos por bucket, pero no equivalen a clique de una instancia arbitraria.\\
& Conjunto independiente & N & Evitar colisiones no se formula como selección de vértices sin adyacencias.\\
& Cobertura de vértices & N & No existe grafo de entrada ni objetivo de cubrir aristas.\\
& Ciclo hamiltoniano & N & La secuencia de buckets no exige visitar vértices una vez.\\
& Errores frecuentes & P & Se reconoce que una analogía estructural no constituye reducción ni equivalencia.\\
& Consecuencias prácticas & P & NP-completitud no implica que problemas polinomiales como esta simulación sean intratables.\\
\midrule
Aproximación & Algoritmos de aproximación & N & El factorial mide políticas; no aproxima una solución óptima NP-difícil.\\
& Razón de aproximación & N & La frontera de Pareto no ofrece razón demostrable respecto de un óptimo único.\\
& Análisis de garantías & P & Se analizan cotas amortizadas/esperadas, que no deben confundirse con garantía de aproximación.\\
\bottomrule
\end{longtable}
\normalsize
